\chapter{Introdução}\label{chp:INTRO}
\section{Apresentação}\label{sec:APRESENT}

Este trabalho apresenta o desenvolvimento de um sistema web destinado à reserva de mesas, salas e auditórios da universidade, permitindo que integrantes da comunidade acadêmica utilizem os espaços de forma organizada e eficiente.

\section{Motivação}\label{sec:MOTIVACAO}
%Explique por que você escolheu esse tema.
%Pode incluir problemas enfrentados atualmente (ex.: conflitos de horários, dificuldades em visualizar os espaços, ocupação inadequada).
%Relacione com benefícios potenciais.

A motivação para este estudo surge da necessidade de otimizar o uso dos espaços da universidade e facilitar o agendamento de ambientes, garantindo que a infraestrutura disponível seja utilizada de forma eficiente e transparente.

\section{Justificativa}\label{sec:Justificativa}
%Argumente por que este trabalho é relevante.
%Pode incluir impacto acadêmico, social ou tecnológico.
%Mencione como a interface baseada nas plantas baixas agrega valor.

O sistema proposto contribui para a gestão inteligente dos espaços universitários, permitindo planejamento e ocupação adequados, além de reduzir conflitos e desperdício de recursos. A utilização das plantas baixas fornecidas pela COPEA garante que a visualização seja precisa, aumentando a confiabilidade do sistema.

\section{Hipótese}\label{sec:HIPOTESE}
%Apresente uma suposição inicial sobre o resultado ou efeito do seu sistema.
%Ex.: o sistema facilitará a reserva de espaços, otimizará ocupação e reduzirá conflitos.
Hipotetiza-se que a implementação de um sistema de reservas com interface de visualização realista permitirá uma gestão mais eficiente dos espaços da universidade, promovendo organização e melhor aproveitamento da infraestrutura disponível.

\section{Objetivos}\label{sec:OBJETIVOS}

\subsection{Objetivos Gerais}\label{sec:OB_GERAIS}
% meta principal do trabalho.
Desenvolver um sistema web para reserva de mesas, salas e auditórios da universidade, com interface de visualização baseada nas plantas baixas, que otimize o uso dos espaços e facilite a organização para a comunidade acadêmica.

\subsection{Objetivos Específicos}\label{sec:OB_Específicos}
% Objetivos Específicos: etapas ou ações para atingir o objetivo geral.
bla bla bla

\begin{itemize}
\item Implementar uma interface gráfica que reproduza fielmente os espaços da universidade;
\item Permitir o agendamento e gerenciamento de reservas por membros da comunidade acadêmica;;
\item Integrar mecanismos de verificação de disponibilidade para evitar conflitos de horários;
\item Garantir usabilidade e acessibilidade na interação do usuário com o sistema.
\end{itemize}

\section{Organização do Trabalho}\label{sec:ORGANIZA}

%Explique como o TCC está estruturado, capítulo por capítulo.
%Curto parágrafo descrevendo o conteúdo de cada capítulo.
O trabalho está organizado da seguinte forma: o Capítulo 1 apresenta a introdução, contextualizando o tema, motivação, justificativa, hipótese e objetivos; o Capítulo 2 aborda a fundamentação teórica e tecnologias utilizadas; o Capítulo 3 descreve o desenvolvimento do sistema, incluindo arquitetura, interface e funcionalidades; o Capítulo 4 apresenta testes, resultados e avaliação do sistema; e, por fim, o Capítulo 5 traz as conclusões, contribuições, limitações e sugestões para trabalhos futuros.

\pagebreak